% Uebungsblatt Allgemeine Chemie (German, university level), with a chemfig drawing.
% Headings: "Aufgabe N (P Punkte)" in bold; parts "a)", "b)", ...
\documentclass[11pt,a4paper]{article}
\usepackage[T1]{fontenc}
\usepackage[utf8]{inputenc}
\usepackage{lmodern}
\usepackage[ngerman]{babel}
\usepackage[margin=2.3cm,top=2.0cm,bottom=2.2cm]{geometry}
\usepackage{amsmath}
\usepackage[version=4]{mhchem}
\usepackage[locale=DE]{siunitx}
\usepackage{chemfig}
\usepackage{caption}
\usepackage{enumitem}

\setlength{\parindent}{0pt}
\setlist[enumerate,1]{label=\alph*), leftmargin=2em, itemsep=0.35em, topsep=0.3em}
\newcommand{\aufgabe}[2]{\bigskip\par\textbf{Aufgabe #1 (#2)}\par\nopagebreak\smallskip}
\setchemfig{atom sep=2.4em}

\begin{document}

% ---- Kopf -----------------------------------------------------------------------------
{\small Allgemeine Chemie \textperiodcentered{} Wintersemester 2026/27 \hfill Institut für Anorganische Chemie}\\[-0.4em]
\rule{\linewidth}{0.6pt}\\[0.5em]
{\Large\bfseries Übungsblatt 5}\\[0.4em]
Ausgabe: 5. November 2026 \hfill Abgabe: Mittwoch, 12. November 2026, 12:00 Uhr\\[0.3em]
{\small Rechenwege müssen nachvollziehbar sein. Verwenden Sie $N_\mathrm{A} = \qty{6.022e23}{mol^{-1}}$.}

% ---- Aufgabe 1 ------------------------------------------------------------------------
\aufgabe{1}{4 Punkte}
Stoffmenge und molare Masse.
\begin{enumerate}
  \item Wie viel Mol sind \qty{36.0}{g} Wasser ($M = \qty{18.0}{g/mol}$), und wie viele Wassermoleküle
        enthält diese Portion?
  \item Welche Masse haben \qty{0.250}{mol} Calciumcarbonat, \ce{CaCO3}? Es gilt
        $A_r(\mathrm{Ca}) = 40{,}08$, $A_r(\mathrm{C}) = 12{,}01$ und $A_r(\mathrm{O}) = 16{,}00$.
  \item Wie viele Wasserstoffatome sind in \qty{1.50}{mol} Methan, \ce{CH4}, enthalten?
\end{enumerate}

% ---- Aufgabe 2 ------------------------------------------------------------------------
\aufgabe{2}{5 Punkte}
Reaktionsgleichungen und Stöchiometrie.
\begin{enumerate}
  \item Im Hochofen reagiert Eisen(III)-oxid mit Kohlenstoffmonoxid zu Eisen und Kohlenstoffdioxid.
        Stellen Sie die ausgeglichene Reaktionsgleichung auf.
  \item Aluminium löst sich in Salzsäure unter Bildung von Aluminiumchlorid und Wasserstoff.
        Gleichen Sie die Gleichung aus: \ce{Al + HCl -> AlCl3 + H2}
  \item Welche Masse Eisen lässt sich aus \qty{160}{g} \ce{Fe2O3} höchstens gewinnen?
        Verwenden Sie $M(\ce{Fe2O3}) = \qty{160}{g/mol}$ und $M(\ce{Fe}) = \qty{56}{g/mol}$.
\end{enumerate}

% ---- Aufgabe 3 ------------------------------------------------------------------------
\aufgabe{3}{6 Punkte}
Ein Acetatpuffer entsteht, wenn \qty{0.10}{mol} Essigsäure ($\mathrm{p}K_\mathrm{s} = 4{,}76$) und
\qty{0.20}{mol} Natriumacetat in Wasser gelöst und auf \qty{1.00}{L} aufgefüllt werden.
\begin{enumerate}
  \item Berechnen Sie den pH-Wert des Puffers mit der Henderson-Hasselbalch-Gleichung.
  \item Zu \qty{1.00}{L} dieses Puffers werden \qty{0.010}{mol} HCl gegeben (Volumenänderung
        vernachlässigen). Welchen pH-Wert hat die Lösung danach?
  \item Erklären Sie in zwei bis drei Sätzen, warum die Mischung aus Essigsäure und Acetat den pH-Wert
        stabilisiert.
  \item Welches Verhältnis $c(\text{Acetat})/c(\text{Essigsäure})$ ist nötig, damit der pH-Wert $4{,}00$
        beträgt?
\end{enumerate}

% ---- Aufgabe 4 ------------------------------------------------------------------------
\aufgabe{4}{5 Punkte}
Lewis-Formel und VSEPR-Modell. Die Abbildung zeigt die Lewis-Formeln von Ammoniak und Bortrifluorid
(freie Elektronenpaare sind als Punktpaare gezeichnet).
\begin{center}
\chemfig{H-\charge{90=\:}{N}(-[6]H)-H}
\hspace{4em}
\chemfig{B(-[:90]\charge{0=\:,180=\:,90=\:}{F})(-[:210]\charge{90=\:,180=\:,270=\:}{F})-[:330]\charge{0=\:,270=\:,90=\:}{F}}
\captionof{figure}{Lewis-Formeln von \ce{NH3} (links) und \ce{BF3} (rechts).}
\end{center}
\begin{enumerate}
  \item Wie viele bindende und wie viele freie Elektronenpaare hat das Zentralatom in \ce{NH3} und in
        \ce{BF3}?
  \item Welche Molekülgeometrie und welche Bindungswinkel sagt das VSEPR-Modell für beide Moleküle voraus?
  \item Warum ist \ce{NH3} ein Dipol, \ce{BF3} aber nicht, obwohl die B--F-Bindung stark polar ist?
  \item \ce{NH3} und \ce{BF3} reagieren zu \ce{H3N-BF3}. Welches Molekül ist die Lewis-Base, welches
        die Lewis-Säure? Begründen Sie kurz.
\end{enumerate}

% ---- Aufgabe 5 ------------------------------------------------------------------------
\aufgabe{5}{4 Punkte}
\qty{5.85}{g} Natriumchlorid ($M = \qty{58.5}{g/mol}$) werden in Wasser gelöst und auf \qty{500}{mL}
aufgefüllt.
\begin{enumerate}
  \item Berechnen Sie die Stoffmengenkonzentration $c$ in mol/L.
  \item Wie groß ist die Massenkonzentration in g/L?
  \item Welches Volumen dieser Lösung enthält \qty{0.0500}{mol} NaCl?
  \item \qty{100}{mL} der Lösung werden mit Wasser auf \qty{400}{mL} verdünnt. Wie groß ist die neue
        Konzentration?
\end{enumerate}

\end{document}
